\documentclass[11pt]{article}
\usepackage[utf8]{inputenc}
\usepackage[margin=1in]{geometry}
\usepackage{amsmath}
\usepackage{graphicx}
\usepackage{booktabs}
\usepackage{hyperref}
\usepackage[round]{natbib}
\usepackage{float}
\usepackage{multirow}
\usepackage{array}

\title{Strong Positive Selection Biases Identity-by-Descent-Based Inferences of Demography and Population Structure in \textit{Plasmodium falciparum}}

\author{%
  Author A$^{1,*}$, Author B$^{2}$, Author C$^{1}$, Author D$^{3}$ \\[6pt]
  $^{1}$Department of Epidemiology, Harvard T.H.\ Chan School of Public Health, Boston, MA, USA \\
  $^{2}$Broad Institute of MIT and Harvard, Cambridge, MA, USA \\
  $^{3}$Faculty of Tropical Medicine, Mahidol University, Bangkok, Thailand \\
  $^{*}$Corresponding author: author.a@hsph.harvard.edu
}
\date{}

\begin{document}
\maketitle

\begin{abstract}
Identity-by-descent (IBD) analysis is widely used for malaria genomic surveillance to estimate effective population size ($N_e$) and infer population structure, yet the impact of strong positive selection from antimalarial drug resistance on these estimates has not been systematically evaluated. We analyzed whole-genome sequencing data from 2,055 \textit{Plasmodium falciparum} isolates from eastern Southeast Asia (SEA; including 640 newly sequenced genomes) spanning 14 years and 18 provinces, with validation in a West African (WAF) dataset. In simulations, strong positive selection shifted IBD segment length distributions toward longer segments ($\Delta$ median length: $+0.42$~cM, Kolmogorov--Smirnov $p < 0.001$), leading to $N_e$ underestimation by 2.3-fold and blurred population structure (adjusted Rand index decreased from 0.91 to 0.54). Empirically, IBD coverage peaks centered on six known drug resistance loci (\textit{kelch13}, \textit{dhfr}, \textit{dhps}, \textit{mdr1}, \textit{crt}, plasmepsin~2/3). An IBD peak removal strategy partially restored accuracy: in WAF, $N_e$ estimates increased 1.8-fold after correction ($p < 0.001$, permutation test), and IBD network community detection recovered geographic structure (adjusted Rand index: 0.52 $\to$ 0.71). In SEA, where high background relatedness dominates (80\% monoclonal, $F_{ws} > 0.95$), correction had minimal effect ($\Delta N_e < 5\%$). These findings demonstrate that selection-aware correction is essential for IBD-based surveillance in high-transmission settings but less necessary where genetic drift predominates.

\medskip
\noindent\textbf{Keywords:} Plasmodium falciparum, identity by descent, selection, drug resistance, population genetics, effective population size, malaria genomic surveillance
\end{abstract}

\section{Introduction}

Malaria genomic surveillance increasingly relies on identity-by-descent (IBD) analysis to estimate effective population size ($N_e$) and infer population structure \citep{taylor2017quantifying, henden2018identity}. IBD segments---contiguous genomic regions inherited from a recent common ancestor---encode information about shared ancestry, migration, and demographic history \citep{browning2018identity, browning2015ibdne}. For \textit{P.~falciparum}, IBD-based approaches have been used to track parasite connectivity, monitor transmission intensity, and detect emerging drug resistance lineages \citep{schaffner2018haplotype}.

However, \textit{P.~falciparum} has undergone some of the strongest selective sweeps observed in any eukaryote, driven by antimalarial drug resistance \citep{ariey2014molecular, miotto2015genetic}. Resistance to chloroquine, sulfadoxine-pyrimethamine, mefloquine, and artemisinin has created extended haplotype blocks around resistance loci, dramatically increasing local IBD sharing \citep{nair2007recurrent, hamilton2017evolution}. These selection-driven IBD increases violate the neutrality assumptions underlying demographic estimators such as IBDNe \citep{browning2015ibdne}, yet no systematic evaluation of this bias has been conducted.

The challenge is compounded by the interplay between selection, genetic drift, and recombination. In low-transmission settings like SEA, high background relatedness from small $N_e$ and population bottlenecks generates long IBD segments genome-wide, potentially masking selection-driven signals \citep{nkhoma2013population}. In high-transmission settings like WAF, low background relatedness makes selection-driven IBD peaks more conspicuous but also more distorting.

Here, we analyze 2,055 \textit{P.~falciparum} genomes from SEA (including 640 newly sequenced) and validate findings in a WAF dataset from MalariaGEN Pf6 \citep{malariagen2021pf7}. Using population genetic simulations and empirical data, we demonstrate that positive selection biases IBD-based inference and develop a correction strategy with context-dependent effectiveness.

\section{Results}

\subsection{Dataset Overview}

Table~\ref{tab:dataset} summarizes the SEA dataset.

\begin{table}[H]
\centering
\caption{Southeast Asian dataset summary. Coverage thresholds indicate the percentage of the genome covered at each depth. $n = 2{,}055$ analyzable isolates.}
\label{tab:dataset}
\begin{tabular}{lc}
\toprule
\textbf{Parameter} & \textbf{Value} \\
\midrule
Total analyzable isolates & 2,055 \\
Newly sequenced & 640 \\
Total in-house & 751 \\
MalariaGEN Pf6 & 1,304 \\
Countries & 4 (Cambodia, Thailand, Laos, Vietnam) \\
Provinces & 18 \\
Time span (years) & 14 \\
$\geq$ 5$\times$ coverage (\% isolates) & 79.3 \\
$\geq$ 10$\times$ coverage (\% isolates) & 68.0 \\
$\geq$ 25$\times$ coverage (\% isolates) & 46.1 \\
\bottomrule
\end{tabular}
\end{table}

\subsection{Monoclonality Assessment}

Table~\ref{tab:monoclonal} compares monoclonal fractions between SEA and WAF.

\begin{table}[H]
\centering
\caption{Monoclonality assessment by within-sample diversity ($F_{ws}$). Monoclonal defined as $F_{ws} > 0.95$. Fisher's exact test for monoclonal proportion comparison.}
\label{tab:monoclonal}
\begin{tabular}{lcccc}
\toprule
\textbf{Population} & \textbf{Monoclonal (\%)} & \textbf{Polyclonal (\%)} & $\boldsymbol{n}$ & $\boldsymbol{p}$ \\
\midrule
SEA & 80.0 & 20.0 & 2,055 & \multirow{2}{*}{$< 0.001$} \\
WAF & 50.7 & 49.3 & 1,246 & \\
\bottomrule
\end{tabular}
\end{table}

\subsection{Simulation: Selection Effects on IBD Distribution}

Simulations with strong positive selection demonstrated three distinct distortions to IBD patterns (Table~\ref{tab:sim_ibd}).

\begin{table}[H]
\centering
\caption{Simulation results: effect of positive selection on IBD metrics. Single-population model with decreasing $N_e$ and high recombination matching \textit{P.~falciparum} parameters. $n = 500$ simulated genome pairs per condition. Kolmogorov--Smirnov test for distribution differences; Cohen's $d$ for effect size.}
\label{tab:sim_ibd}
\begin{tabular}{lcccc}
\toprule
\textbf{IBD Metric} & \textbf{Neutral} & \textbf{Under Selection} & $\boldsymbol{p}$ (KS) & \textbf{Cohen's $d$} \\
\midrule
Median segment length (cM) & $1.24 \pm 0.31$ & $1.66 \pm 0.48$ & $< 0.001$ & 1.04 \\
Mean total pairwise IBD (cM) & $3.87 \pm 1.42$ & $6.21 \pm 2.18$ & $< 0.001$ & 1.27 \\
Chromosomal IBD uniformity (CV) & 0.18 & 0.67 & $< 0.001$ & --- \\
\bottomrule
\end{tabular}
\end{table}

\subsection{Simulation: Effects on $N_e$ Estimation}

Selection-induced IBD inflation caused IBDNe to underestimate $N_e$ (Table~\ref{tab:sim_ne}).

\begin{table}[H]
\centering
\caption{Effect of selection on $N_e$ estimation via IBDNe. True $N_e$ set to 10,000 in simulations. Bias = estimated/true $N_e$. $n = 100$ simulation replicates per condition. Wilcoxon signed-rank test for bias $\neq 1$.}
\label{tab:sim_ne}
\begin{tabular}{lccc}
\toprule
\textbf{Condition} & $\boldsymbol{\hat{N}_e}$ \textbf{(M $\pm$ SD)} & \textbf{Bias (fold)} & $\boldsymbol{p}$ \\
\midrule
Neutral & $9,840 \pm 1,420$ & 0.98 & 0.61 \\
Under selection & $4,350 \pm 980$ & 0.44 & $< 0.001$ \\
After peak removal & $7,210 \pm 1,650$ & 0.72 & 0.004 \\
\bottomrule
\end{tabular}
\end{table}

\subsection{Simulation: Effects on Population Structure}

In a five-deme stepping-stone model, selection blurred population structure (Table~\ref{tab:sim_structure}).

\begin{table}[H]
\centering
\caption{Stepping-stone simulation: effect of selection on population structure inference. Five subpopulations with symmetric migration. Adjusted Rand index (ARI) measures agreement between inferred and true community assignments (1 = perfect). $n = 50$ simulation replicates.}
\label{tab:sim_structure}
\begin{tabular}{lccc}
\toprule
\textbf{Condition} & \textbf{ARI (M $\pm$ SD)} & \textbf{Communities Detected} & $\boldsymbol{p}$ \textbf{vs.\ Neutral} \\
\midrule
Neutral & $0.91 \pm 0.05$ & 5.0 & --- \\
Under selection & $0.54 \pm 0.12$ & 3.2 & $< 0.001$ \\
After peak removal & $0.73 \pm 0.09$ & 4.1 & $< 0.001$ \\
\bottomrule
\end{tabular}
\end{table}

\subsection{Empirical IBD Peaks at Drug Resistance Loci}

IBD coverage profiles showed peaks centered on six known drug resistance genes (Table~\ref{tab:peaks}).

\begin{table}[H]
\centering
\caption{Drug resistance loci identified in IBD coverage peaks. Peak height = fold-enrichment of IBD proportion at locus vs.\ genome-wide median. $n = 2{,}055$ SEA isolates.}
\label{tab:peaks}
\begin{tabular}{llccc}
\toprule
\textbf{Gene} & \textbf{Drug Resistance} & \textbf{Chromosome} & \textbf{Peak Height (fold)} & \textbf{Peak Width (kb)} \\
\midrule
\textit{kelch13} & Artemisinin & 13 & 8.4 & 320 \\
\textit{dhfr} & Pyrimethamine & 4 & 6.2 & 180 \\
\textit{dhps} & Sulfadoxine & 8 & 5.7 & 210 \\
\textit{mdr1} & Mefloquine/chloroquine & 5 & 4.9 & 250 \\
\textit{crt} & Chloroquine & 7 & 4.3 & 190 \\
Plasmepsin 2/3 & Piperaquine & 14 & 7.1 & 280 \\
\bottomrule
\end{tabular}
\end{table}

\subsection{SEA vs.\ WAF: Context-Dependent Correction Effectiveness}

Table~\ref{tab:correction} compares the effect of IBD peak removal between SEA and WAF.

\begin{table}[H]
\centering
\caption{Context-dependent effectiveness of IBD peak removal. $\Delta N_e$ = change in estimated $N_e$ after peak removal. $\Delta$ARI = change in adjusted Rand index for community detection. Permutation test ($n = 10{,}000$) for $N_e$ change significance.}
\label{tab:correction}
\begin{tabular}{lcccccc}
\toprule
\textbf{Feature} & \multicolumn{2}{c}{\textbf{SEA}} & \multicolumn{2}{c}{\textbf{WAF}} & $\boldsymbol{p}$ \textbf{(SEA)} & $\boldsymbol{p}$ \textbf{(WAF)} \\
\cmidrule(lr){2-3} \cmidrule(lr){4-5}
& \textbf{Before} & \textbf{After} & \textbf{Before} & \textbf{After} & & \\
\midrule
$\hat{N}_e$ (relative) & 1.00 & 1.04 & 1.00 & 1.82 & 0.38 & $< 0.001$ \\
ARI (community) & 0.78 & 0.80 & 0.52 & 0.71 & 0.41 & $< 0.001$ \\
Monoclonal fraction & \multicolumn{2}{c}{80\%} & \multicolumn{2}{c}{50.7\%} & --- & --- \\
Background relatedness & \multicolumn{2}{c}{High} & \multicolumn{2}{c}{Low} & --- & --- \\
\bottomrule
\end{tabular}
\end{table}

\section{Discussion}

This study demonstrates that strong positive selection from antimalarial drug resistance systematically biases IBD-based estimates of $N_e$ and population structure in \textit{P.~falciparum}. The bias operates through a clear mechanism: selection creates extended haplotype blocks that inflate IBD segment lengths and total pairwise IBD sharing. IBDNe interprets these longer segments as evidence for smaller or more recent $N_e$, while community detection algorithms merge truly distinct subpopulations that share resistance haplotypes.

The context dependence of our correction strategy has practical implications. In high-transmission settings like WAF, where background relatedness is low and selection-driven IBD peaks are conspicuous against the genomic background, peak removal substantially improves both $N_e$ estimation and structure inference. In low-transmission settings like SEA, where genetic drift has independently inflated background IBD sharing, selection-driven peaks contribute relatively less to the total IBD signal, and correction has minimal effect.

These findings have immediate relevance for malaria elimination programs that use genomic surveillance to monitor transmission intensity and parasite connectivity \citep{who2022malaria}. Programs relying on IBD-based $N_e$ estimates to track transmission trends should account for selection bias, particularly in Africa where the discrepancy between corrected and uncorrected estimates is substantial.

Limitations include the reliance on known drug resistance loci for peak identification; novel or uncharacterized selected loci would be missed. The peak removal strategy is conservative and may not eliminate all selection-driven IBD inflation if selection signals extend beyond defined peak boundaries. Future work should develop model-based approaches that jointly estimate selection and demography from IBD data.

\section{Methods}

\subsection{Samples and Sequencing}
WGS data from 2,055 eastern SEA isolates (751 sequenced in-house including 640 new; 1,304 from MalariaGEN Pf6 \citep{malariagen2021pf7}) spanning Cambodia, Thailand, Laos, and Vietnam. WAF validation from MalariaGEN Pf6.

\subsection{Monoclonality Assessment}
Within-sample diversity ($F_{ws}$) \citep{manske2012analysis} was used to classify isolates as monoclonal ($F_{ws} > 0.95$) or polyclonal. Polyclonal infections were deconvolved with dEploid \citep{zhu2017deploid}.

\subsection{IBD Inference}
IBD segments were inferred using haplotype-based methods \citep{browning2018identity}. $N_e$ was estimated using IBDNe \citep{browning2015ibdne}. Population structure was inferred via Louvain community detection \citep{blondel2008fast} on IBD networks.

\subsection{Simulations}
Single-population and five-deme stepping-stone models with strong positive selection, decreasing $N_e$, and high recombination matching \textit{P.~falciparum} parameters. True IBD from simulation pedigrees used for validation.

\subsection{IBD Peak Removal}
Genomic regions with excess IBD sharing ($>$3 SD above genome-wide mean) were identified and excluded before downstream analysis.

\section{Conclusion}

Strong positive selection from drug resistance biases IBD-based demographic and structural inference in \textit{P.~falciparum}. The practical impact of this bias, and the effectiveness of correction, depend critically on the epidemiological setting: correction is essential in high-transmission regions with low background relatedness but largely unnecessary where genetic drift dominates. These findings should guide the design and interpretation of malaria genomic surveillance programs worldwide.

\bibliographystyle{plainnat}
\bibliography{references}

\end{document}
